\begin{Theorem}
\label{thm:compact-perturbation-fredholm-index-invariance}
Let $A\in\mathcal{L}(X,Y)$ be semi-Fredholm and let $k\in\mathcal{L}(X,Y)$ be compact. Then $A+k$ is semi-Fredholm and $\operatorname{ind}A+k=\operatorname{ind}A$.
\end{Theorem}

\begin{Proof}
We first consider the special case where $Y=X$ and $A=I$. We have to show that $I+k$ is Fredholm with index zero. Let $x_n\in\mathcal{N}(I+k)$ with $\|x_n\|\leq1$. After taking a subsequence, we may assume that $kx_n$ converges, and since $x_n+kx_n=0$, we conclude that $x_n$ converges. Hence the unit ball in $\mathcal{N}(I+k)$ is compact, which implies that $\mathcal{N}(I+k)$ is finite-dimensional.

We next show that $\mathcal{R}(I+k)$ is closed. Let $V$ be a complement of $\mathcal{N}(I+k)$, then $I+k$ is a bijection from $V$ to $\mathcal{R}(I+k)$. If we show that the inverse is bounded, then $\mathcal{R}(I+k)$ is isomorphic to $V$, hence complete. Suppose now that the inverse is not bounded, i.e., there is a sequence $x_n\in V$ such that $\|x_n\|=1$, but $x_n+kx_n\to0$. Then we may again take a subsequence and assume that $kx_n$ converges, say to $y$. It follows that $x_n\to-y$ and that $y+ky=0$, i.e., $y\in\mathcal{N}(I+k)$. But this is a contradiction, since on the other hand, $y\in V$ and $\|y\|=1$.

We have thus proved that $I+k$ is semi-Fredholm. That the index is zero follows by considering the family $\{I+tk\mid t\in[0,1]\}$.

For the general case, we shall again assume that either $A$ is Fredholm or that $X$ and $Y$ are Hilbert spaces. Let $X=V\oplus\mathcal{N}(A)$, $Y=W\oplus\mathcal{R}(A)$. Let $S\in\mathcal{L}(Y,X)$ be defined as follows: If $y\in\mathcal{R}(A)$, then $Sy$ is the unique $x\in V$ such that $Ax=y$, and if $y\in W$, then $Sy=0$. If $W$ is finite-dimensional, then $AS$ is a finite rank perturbation of the identity, i.e., $AS-I$ is compact. Hence $(A+k)S-I$ is also compact, i.e., $(A+k)S$ is Fredholm with index zero. Hence $\mathcal{R}((A+k)S)$ and consequently $\mathcal{R}(A+k)$ has finite codimension.

If on the other hand $\mathcal{N}(A)$ is finite-dimensional, then $SA-I$ is a finite rank operator and $S(A+k)$ is Fredholm with index zero. Hence the nullspace of $S(A+k)$ and a fortiori the nullspace of $A+k$ are finite-dimensional. We need to show that the range of $A+k$ is closed. Let $U$ be a complement in $X$ to the linear span of $\mathcal{N}(A)$ and $\mathcal{N}(A+k)$. We shall show that $A+k|_U$ has a bounded inverse. This implies that $(A+k)(U)$ is complete, and since this space has finite codimension in $\mathcal{R}(A+k)$, the closedness of $\mathcal{R}(A+k)$ follows. Assume now that $x_n\in U$ with $\|x_n\|=1$ and $(A+k)x_n\to0$. After taking a subsequence, we may assume that $kx_n$ and hence $Ax_n$ converges. But $A|_U$ has a bounded inverse, hence $x_n$ converges, say to $y$. As above, we find a contradiction, since $(A+k)y=0$, but on the other hand $y\in U$ with $\|y\|=1$.

Hence we have shown that $A+k$ is semi-Fredholm and the statement about the index follows again by considering the family $\{A+tk\mid t\in[0,1]\}$.
\end{Proof}

Now let $A$ be a densely defined, closed symmetric operator in a Hilbert space $H$. If $\lambda\in\mathbb{C}\backslash\mathbb{R}$, then $A-\lambda I$ has a bounded inverse; hence it is injective and has closed range. By Corollary~\ref{cor:semi-fredholm-open-set-constant-index}, the deficiency of $A-\lambda I$ is constant in the upper and lower half-planes. Thus the Fredholm index of $A-\lambda I$ for $\lambda$ in the upper half-plane is equal to minus the deficiency index $\gamma^+$ of $A$, and the Fredholm index in the lower half-plane is equal to $-\gamma^-$ (cf. Definition~\ref{def:operator-deficiency-indices}). By Theorem~\ref{thm:essential-self-adjointness-zero-deficiency}, $A$ is self-adjoint if and only if both deficiency indices are zero.

\Needspace{5\baselineskip}\textbf{Problems}

\begin{enumerate}[label=\thesection.\arabic*.,ref=\thesection.\arabic*,start=54,series=book-problems]

\item\label{ex:closed-plus-finite-dimensional-subspace} Let $V$ be a closed subspace of $X$ and let $W$ be a finite-dimensional subspace. Prove that $V+W$ is closed in $X$.

\item\label{ex:positive-potential-dirichlet-classical-solvability} Let $c\in C([0,1])$ be such that $c(x)\geq0$. Prove that the equation $y\prime\prime-c(x)y=f(x)$ subject to boundary conditions $y(0)=y(1)=0$ has a solution in $C^2([0,1])$ for every given $f\in C([0,1])$.

\item\label{ex:half-line-differentiation-deficiency-indices} Let $H=L^2(0,\infty)$ and let $Au=iu\prime$ with domain $H_0^1(0,\infty)$. Find the deficiency indices of $A$ and its adjoint.

\item\label{ex:real-resolvent-implies-self-adjointness} Let $A$ be closed, densely defined and symmetric and such that the resolvent set of $A$ contains at least one real number. Prove that $A$ is self-adjoint. If in addition $A$ is positive definite, show that the negative real axis belongs to the resolvent set of $A$.

\end{enumerate}
